In the second part of the definition, one has therefore used identifications
and abuses of notation in current
use%
{\renewcommand{\thefootnote}{(8)}\nde{: for example, one has identified,
on the one hand,
\(\mathrm{pr}_{2,1}^{*}(X''_{1})=\mathrm{pr}_{3,1}^{*}(X''_{1})\) and, on the
other hand,
\(\mathrm{pr}_{2,1}^{*}(X''_{2})=\mathrm{pr}_{3,2}^{*}(X''_{2})\).}},
which experience proves to be harmless, but which it is obviously advisable
to avoid in a rigorous expos\'e of the theory of descent, (which must
precisely justify to a certain extent these current abuses of language).
Such a formal expos\'e ([D]) has been written by Giraud, (with a view to
justifying and making precise SGA~1, VII, which has never been written).
For a detailed expos\'e of the results of faithfully flat descent, of which
constant use will be made in the present Seminar, one may consult
SGA~1, VIII.

Let \(f\colon S'\to S\) still be a morphism such that \(S''=S'\times_{S}S'\)
exists, and let \(X\) be an object over \(S\) such that \(X'=X\times_{S}S'\)
and \(X''=X\times_{S}S''\) exist; then the inverse images of \(X'\) by
\(\mathrm{pr}_{i}\) (\(i=1,2\)) exist and are canonically isomorphic, and
consequently \(X'/S'\) is endowed with a canonical gluing datum relative
to~\(f\). When \(S'''\) and \(X'''=X\times_{S}S'''\) exist, it is even a
descent datum. If \(Y\) is another object over \(S\), satisfying the same
conditions as \(X/S\), then for every \(S\)-morphism \(X\to Y\), the
corresponding \(S'\)-morphism \(X'\to Y'\) is ``compatible with the
canonical gluing data'' on \(X'\), \(Y'\). If in particular \(S'\to S\) is a
\emph{quarrable} morphism, then
\[
X\longmapsto X'=X\times_{S}S'
\]
is a functor from the category \(\mathcal{C}_{/S}\) into the ``category of
objects over \(S'\) endowed with a descent datum relative to~\(f\)''\,---\,a
category whose definition is left to the reader, and which is a full
subcategory of the ``category of objects over \(S'\) endowed with a gluing
datum relative to~\(f\)''.

This being said:

\begin{definition}{2.2}
\label{IV.2.2}
One says that a morphism \(f\colon S'\to S\) is a \emph{descent morphism}
(\resp an \emph{effective descent morphism}) if \(f\) is \emph{quarrable}
(\ie for every \(X\to S\), the fiber product \(X\times_{S}S'\) exists), and
if the preceding functor \(X\mapsto X'=X\times_{S}S'\) from the category
\(\mathcal{C}_{/S}\) of objects over \(S\), into the category of objects
over \(S'\) endowed with a descent datum relative to~\(f\), is fully faithful
(\resp an equivalence of categories).
\end{definition}

One will note that the first of these two notions can be expressed with the
help of only the notion of gluing datum (hence without involving the triple
fiber product \(S'''\)), \(f\) being a descent morphism if \(f\) is
\emph{quarrable}
and \(X\mapsto X'\) is a fully faithful functor from the category
\(\mathcal{C}_{/S}\) into the category of objects over \(S'\) endowed with a
gluing datum relative to~\(f\). When one makes this definition explicit, one
finds that this means that for two objects \(X\), \(Y\) over \(S\), the
following diagram of sets
\[
(\mathrm{x})\qquad
\xymatrix@C=2.4em{
\Hom_{S}(X,Y) \ar[r]
  & \Hom_{S'}(X',Y')
      \ar@<0.55ex>[r]^{p_{1}}
      \ar@<-0.55ex>[r]_{p_{2}}
  & \Hom_{S''}(X'',Y'')
}
\]
is exact, where \(p_{i}(u')\) denotes the inverse image of
\(u'\in\Hom_{S'}(X',Y')\) by the projection
\(\mathrm{pr}_{i}\colon S''\to S'\), for \(i=1,2\); indeed, the kernel of
the pair \((p_{1},p_{2})\) is by definition nothing other than the set of
\(S'\)-morphisms \(X'\to Y'\) compatible with the canonical gluing data.

Let us note that, by definition of the inverse images \(Y'\), \(Y''\), one
has canonical bijections
\[
\Hom_{S'}(X',Y')\simeq\Hom_{S}(X',Y)
\qquad\text{and}\qquad
\Hom_{S''}(X'',Y'')\simeq\Hom_{S}(X'',Y),
\]
so that the exactness of the diagram~\((\mathrm{x})\) is equivalent to that
of
\[
(\mathrm{xx})\qquad
\xymatrix@C=2.8em{
\Hom_{S}(X,Y) \ar[r]
  & \Hom_{S}(X',Y)
      \ar@<0.55ex>[r]
      \ar@<-0.55ex>[r]
  & \Hom_{S}(X'',Y)\,,
}
\]
obtained by applying \(\Hom_{S}(-,Y)\) to the diagram in
\(\mathcal{C}_{/S}\):
\[
(\mathrm{xxx})\qquad
\xymatrix@C=3em{
X'' \ar@<0.5ex>[r] \ar@<-0.5ex>[r] & X' \ar[r] & X
}
\]
which is deduced from
\[
\xymatrix@C=3em{
S'' \ar@<0.5ex>[r] \ar@<-0.5ex>[r] & S' \ar[r] & S
}
\]
by the change of base \(X\to S\). This proves, taking~1.6 into account, the
first part of the following proposition:

\begin{proposition}{2.3}
\label{IV.2.3}
Let \(f\colon S'\to S\) be a morphism. If it is a universal effective
epimorphism, it is a descent morphism, and the converse is true if \(S\) is
a quarrable object of \(\mathcal{C}\) (\ie its product by every object of
\(\mathcal{C}\) exists).
\end{proposition}

It remains to prove that if \(f\) is a descent morphism, it is a universal
effective epimorphism, that is, that for every morphism \(X\to S\), the
diagram~\((\mathrm{xxx})\) is exact, \ie for every object \(Z\) of
\(\mathcal{C}\), the transform of this diagram by \(\Hom(-,Z)\) is an exact
diagram of sets. Now by hypothesis \(Z\times S\) exists; let \(Y\) be the
object of \(\mathcal{C}_{/S}\) which it defines; then the transform of
diagram~\((\mathrm{xxx})\) by \(\Hom(-,Z)\) is isomorphic to the transform
by \(\Hom_{S}(-,Y)\), now the latter is exact by the hypothesis on~\(f\).

One may therefore apply to universal effective epimorphisms the results on
descent morphisms, such as the following:

\begin{proposition}{2.4}
\label{IV.2.4}
Let \(f\colon S'\to S\) be a descent morphism (for example a universal
effective epimorphism). Then:

a)~For every \(S\)-morphism \(u\colon X\to Y\), \(u\) is an isomorphism
(\resp a monomorphism) if and only if \(u'\colon X'\to Y'\) is.

b)~Let \(X\), \(Y\) be two subobjects of \(S\), and \(X'\) and \(Y'\) the
subobjects of \(S'\), inverse images of the preceding ones. In order that
\(X\) be majorized by \(Y\) (\resp be equal to \(Y\)), it is necessary and
sufficient that the same hold for \(X'\), \(Y'\).
\end{proposition}

For~(a), it follows from the definition that if \(u'\) is an isomorphism in
the category of objects with gluing datum, then \(u\) is an isomorphism; now
one sees at once that every isomorphism of objects over \(S'\), compatible
with gluing data, is an isomorphism of objects with gluing datum,
% typo: un isomorphismes -> un isomorphisme
\ie its inverse is likewise compatible with the gluing data. For~b), one is
reduced to proving that if \(X'\) is majorized by \(Y'\), \ie if there
exists an \(S'\)-morphism \(X'\to Y'\), then the same holds for \(X\), \(Y\)
over \(S\). Now since \(Y'\to S'\), hence also \(Y''\to S''\), is a
monomorphism, one sees that \(X'\to Y'\) is automatically compatible with
the gluing data, hence comes from an \(S\)-morphism \(X\to Y\). Let us note
that the proof holds more generally when one has two objects \(X\), \(Y\)
over \(S\), with \(Y\to S\) a monomorphism, and one asks whether the morphism
\(X\to S\) factors through \(Y\): it suffices that \(X'\to S'\) factor
through \(Y'\).

\begin{corollary}{2.5}
\label{IV.2.5}
Let \(f\colon S'\to S\) be a universal effective epimorphism and
\(g\colon S\to T\) a morphism such that \(S\times_{T}S\) exists. Suppose that
\(S''=S'\times_{S}S'\) is also a fiber product of \(S'\) by itself over \(T\),
\ie \(S'\times_{S}S'\isomto S'\times_{T}S'\). Then \(g\colon S\to T\) is a
monomorphism (and conversely, of course).
\end{corollary}

Indeed, let us consider the cartesian diagram
\[
\xymatrix@C=4.2em@R=2.4em{
S'\times_{S}S' \ar[r]^{\sim} \ar[d]
  & S'\times_{T}S' \ar[d]^{f} \\
S \ar[r]
  & S\times_{T}S\,,
}
\]
where the second horizontal arrow is the diagonal morphism. By virtue
of~1.9 the second vertical arrow \(f\) is a universal effective
epimorphism, by hypothesis the first horizontal arrow is an isomorphism,
hence by virtue of~2.4~a) or~b), at one's
choice%
{\renewcommand{\thefootnote}{(9)}\nde{: applied to
\(f\colon S'\times_{T}S'\to S\times_{T}S=Y\).}},
the same holds of \(S\to S\times_{T}S\), which means precisely that
\(g\colon S\to T\) is a monomorphism.

\begin{remark}{2.6}
\label{IV.2.6}
The notions introduced in~2.1, in terms of morphisms between certain
projective limits, are made explicit in an evident way in terms of the
contravariant functors defined by the objects \(S\), \(S'\), \(X'\)
envisaged: subject to the existence of the fiber products envisaged
in~2.1, there is a one-to-one correspondence between the gluing data
(\resp descent data) on \(X'/S'\) relative to \(S'\to S\), and the gluing
data (\resp descent data) for the corresponding objects in
\(\widehat{\mathcal{C}}=\Hom(\mathcal{C}^{\circ},\Ens)\). This makes it
possible, if one so desires, to extend these notions to the case where one
makes no existence hypothesis on fiber products in \(\mathcal{C}\).
\end{remark}

\begin{remark}{2.7}
\label{IV.2.7}
The notions introduced in this number generalize to the case of families of
morphisms. They can on the other hand be presented in a more intrinsic
manner with the help of the notion of sieve~(4.1). For these questions, the
reader is referred to~[D].
\end{remark}

\sgasection{3}{Universal effective equivalence relations}
\label{IV.sec.3}

\par\addvspace{0.55\baselineskip}
\noindent\textbf{3.1. Equivalence relations: definitions}\label{IV.sec.3.1}
\par\addvspace{0.35\baselineskip}

\begin{definition}{3.1.1}
\label{IV.3.1.1}
A \emph{representable subfunctor} \(R\) of the functor \(X\times X\), such
that for every \(S\in\Ob(\mathcal{C})\), \(R(S)\) is the graph of an
equivalence relation in \(X(S)\), is called a
\emph{\(\mathcal{C}\)-equivalence relation} in \(X\in\Ob(\mathcal{C})\).
\end{definition}

This definition applies in particular to the category
\(\widehat{\mathcal{C}}\). If one considers \(X\) as an object of
\(\widehat{\mathcal{C}}\), one sees then that a
\(\widehat{\mathcal{C}}\)-equivalence relation in \(X\) is nothing other
than a subfunctor \(R\) of \(X\times X\) such that \(R(S)\) is the graph of
an equivalence relation in \(X(S)\) for every object \(S\) of
\(\mathcal{C}\).%
{\renewcommand{\thefootnote}{(10)}\nde{: The condition is obviously
necessary. Conversely, if for every \(S\in\Ob(\mathcal{C})\), \(R(S)\) is
the graph of an equivalence relation, then this equivalence relation extends
to \(R(F)\) for every \(F\in\Ob(\widehat{\mathcal{C}})\), by declaring that
two morphisms \(\varphi,\psi\colon F\to R\) are equivalent if, for every
\(S\in\Ob(\mathcal{C})\) and \(x\in F(S)\), \(\varphi(x)\) and \(\psi(x)\) are
equivalent in \(X(S)\).}}
% typo?: footnote (10) prints \(\varphi,\psi\colon F\to R\); kept as printed

If \(R\) is a \(\mathcal{C}\)-equivalence relation in \(X\), one denotes by
\(p_{i}\) the morphism from \(R\) into \(X\) induced by the projection
\(\mathrm{pr}_{i}\colon X\times X\to X\). One therefore has a diagram
\[
p_{1},p_{2}\colon R\rightrightarrows X.
\]

\begin{definition}{3.1.2}
\label{IV.3.1.2}
A morphism \(u\colon X\to Z\) is said to be \emph{compatible} with \(R\) if
\(up_{1}=up_{2}\). An object of \(\mathcal{C}\), cokernel of the pair
\((p_{1},p_{2})\), is also called a \emph{quotient-object} of \(X\) by \(R\)
and denoted \(X/R\).
\end{definition}

One therefore has an exact diagram
\[
\xymatrix@C=3.6em{
R \ar@<0.55ex>[r]^{p_{1},p_{2}} \ar@<-0.55ex>[r]
  & X \ar[r]^{p}
  & X/R
}
\]
and \(X/R\) represents the \emph{covariant} functor
\[
\Hom_{\mathcal{C}}(X/R,Z)
=\{\text{morphisms from \(X\) into \(Z\) compatible with \(R\)}\}.
\]
As the quotient-objects have been chosen in \(\mathcal{C}\), the quotient
\(X/R\) is unique (when it exists).

These definitions generalize at once to the case of a
\(\widehat{\mathcal{C}}\)-equivalence relation in \(X\), but one will
observe that the identification of the objects of \(\mathcal{C}\) with
their images in \(\widehat{\mathcal{C}}\) \emph{does not commute with the
formation of quotients}, that is, that the quotient \(X/R\) of \(X\) by
\(R\) in \(\mathcal{C}\) is not \emph{a priori} a quotient of \(X\) by \(R\)
in \(\widehat{\mathcal{C}}\). One will therefore take care not to identify
\(\mathcal{C}\) rashly with its image in \(\widehat{\mathcal{C}}\) in
questions involving passages to the
quotient.%
{\renewcommand{\thefootnote}{(11)}\nde{: Let us illustrate this remark by
giving a glimpse of the sequel of this Expos\'e. Let \(G\) be a
\(\mathcal{C}\)-group, \(H\) a \(\mathcal{C}\)-subgroup, \(R\) the
equivalence relation in \(G\) defined by \(G\times H\to G\times G\),
\((g,h)\mapsto(g,gh)\) (\cf~3.2). The functor \(Q\) defined by
\(Q(S)=G(S)/H(S)\) is a quotient in \(\widehat{\mathcal{C}}\) (by~4.4.9
applied to the coarsest topology, \cf~4.4.2), but it is not in general the
quotient that one wishes. For example, for \(\mathcal{C}=\Sch\), one has an
exact sequence of (affine) group schemes:
\[1\longrightarrow\bmu_{2}\longrightarrow\Gm\xrightarrow{p}\Gm\longrightarrow1\]
which identifies \(\Gm\) with the quotient \(\Gm/\bmu_{2}\). Moreover, as
\(p\) is a finite and locally free morphism, \(\Gm\) is then the
\emph{quotient-sheaf} of \(\Gm\) by \(\bmu_{2}\) in the larger category of
\emph{sheaves for the \textup{(fppf)} topology}, \cf~4.6.6\,(ii) and~6.3.2. By
contrast, the quotient \(Q\) in \(\widehat{\mathcal{C}}\) is not isomorphic
to \(\Gm\) since, for example, \(Q(\ZZ)=\{1\}\) whereas
\(\Gm(\ZZ)=\{\pm1\}\). Hence \(Q\) is not an \textup{(fppf)} sheaf, and a
fortiori \(Q\) is not representable.}}

In the sequel, \emph{we shall say simply equivalence relation for
\(\widehat{\mathcal{C}}\)-equivalence relation}; we shall specify, when the
case arises, whether it is a matter of \(\mathcal{C}\)-equivalence
relations.%
% typo: relations d'équivalences -> relations d'équivalence
{\renewcommand{\thefootnote}{(12)}\nde{: Take care that, even for a
\(\widehat{\mathcal{C}}\)-equivalence relation in \(X\), one is interested
in the existence of a quotient \emph{in \(\mathcal{C}\)}. }}

\begin{definition}{3.1.3}
\label{IV.3.1.3}
If \(X\) is an object of \(\mathcal{C}\) over \(S\), an equivalence relation
\(R\) in \(X\) such that the structural morphism \(X\to S\) is compatible
with \(R\) is called an equivalence relation in \(X\) over \(S\).
\end{definition}

The canonical monomorphism \(R\to X\times X\) then factors through the
monomorphism
\[
X\times_{S}X\longrightarrow X\times X
\]
and defines an equivalence relation in the object \(X\to S\) of
\(\mathcal{C}_{/S}\). When the quotient \(X/R\) exists, it is endowed with
a canonical morphism into \(S\) and the corresponding object of
\(\mathcal{C}_{/S}\) is a quotient of \(X\in\Ob(\mathcal{C}_{/S})\) by the
preceding equivalence relation. Conversely, if \(S\) is a \emph{quarrable}
object of \(\mathcal{C}\) and if \(Y\to S\) is a quotient of \(X\) by this
equivalence relation (in \(\mathcal{C}_{/S}\)), then \(Y\) is a quotient of
\(X\) by \(R\) in \(\mathcal{C}\). In any case, we shall never have to
consider quotients in \(\mathcal{C}_{/S}\) which are not already quotients
in \(\mathcal{C}\).

\begin{definition}{3.1.4}
\label{IV.3.1.4}
If \(X\) (\resp\ \(X'\)) is an object of \(\mathcal{C}\) endowed with an
equivalence relation \(R\) (\resp \(R'\)), a morphism
\[
u\colon X\to X'
\]
is said to be \emph{compatible with \(R\) and \(R'\)} if the following
equivalent conditions are satisfied:
\begin{enumeratei}
\item for every \(S\in\Ob(\mathcal{C})\), two points of \(X(S)\) congruent
modulo \(R(S)\) are transformed by \(u\) into two points of \(X'(S)\)
congruent modulo \(R'(S)\);
\item there exists a morphism \(R\to R'\) (necessarily unique) making the
following diagram commutative
\[
\xymatrix@C=4.5em@R=2.6em{
R \ar[r] \ar[d]
  & R' \ar[d] \\
X\times X \ar[r]_{u\times u}
  & X'\times X'\,.
}
\]
\end{enumeratei}
\end{definition}

By the universal property of \(X/R\), there then exists (when the quotients
\(X/R\) and \(X'/R'\) exist) a unique morphism \(v\) making the following
diagram commutative
\[
\xymatrix@C=4em@R=2.4em{
X \ar[r]^{p} \ar[d]_{u}
  & X/R \ar[d]^{v} \\
X' \ar[r]_{p'}
  & X'/R'\,.
}
\]

\begin{definition}{3.1.5}
\label{IV.3.1.5}
A subobject (\(={}\) a representable subfunctor) \(Y\) of \(X\) is said to
be \emph{stable} under the equivalence relation \(R\) if the following
equivalent conditions are satisfied:
\begin{enumeratei}
\item For every \(S\in\Ob(\mathcal{C})\), the subset \(Y(S)\) of \(X(S)\) is
stable under \(R(S)\).
\item The two subobjects of \(R\), inverse images of \(Y\) by \(p_{1}\) and
\(p_{2}\), are identical.
\end{enumeratei}
\end{definition}

An important particular case of a stable subobject of \(X\) is the
following: the quotient \(X/R\) exists and \(Y\) is the inverse image on
\(X\) of a subobject of \(X/R\).

\begin{definition}{3.1.6}
\label{IV.3.1.6}
Let \(R\) be an equivalence relation in \(X\) and \(X'\to X\) a morphism.
The equivalence relation \(R'\) in \(X'\) obtained by the cartesian diagram
\[
\xymatrix@C=4.5em@R=2.6em{
R' \ar[r] \ar[d]
  & R \ar[d] \\
X'\times X' \ar[r]
  & X\times X
}
\]
is called the equivalence relation in \(X'\), \emph{inverse image} of the
equivalence relation \(R\) in \(X\). In particular, if \(X'\) is a subobject
of \(X\), one will say that it is the equivalence relation \emph{induced}
in \(X'\) by \(R\); one will denote it by \(R_{X'}\).
\end{definition}

The morphism \(X'\to X\) is compatible with \(R'\) and \(R\); one therefore
has, when the quotients exist, a morphism \(X'/R'\to X/R\) (3.1.4). If
\(X'\) is a subobject of \(X\), we shall see later that in certain cases one
can prove that \(X'/R'\to X/R\) is a monomorphism, hence identifies
\(X'/R'\) with a subobject of \(X/R\). When this is the case, the inverse
image of this subobject in \(X\) will be a subobject of \(X\) majorizing
\(X'\) and stable under \(R\): the \emph{saturation} of \(X'\) for the
equivalence relation \(R\).

\begin{proposition}{3.1.7}
\label{IV.3.1.7}
If the subobject \(Y\) of \(X\) is stable under \(R\), one has two cartesian
squares, for \(i=1,2\):
\[
\xymatrix@C=4em@R=2.4em{
R_{Y} \ar[r] \ar[d]_{p_{i}}
  & R \ar[d]^{p_{i}} \\
Y \ar[r]
  & X\,.
}
\]
\end{proposition}

\begin{proof}
Immediate.
\end{proof}
